\documentclass[11pt, twocolumn, landscape, a4paper]{article}

\usepackage[margin=1cm]{geometry}
\usepackage[utf8]{inputenc}
\usepackage[spanish]{babel}
\usepackage{amssymb, amsmath, amsbsy}
\usepackage{tasks}
\usepackage{enumerate}
\usepackage{enumitem}
\usepackage{graphicx}

\usepackage{tikz}
\usetikzlibrary{angles,quotes}

\usepackage[pdftex]{hyperref}
\hypersetup{pdftitle={Trigonometría},pdfauthor={Luis Carrillo Gutiérrez}}

\fontfamily{cmss}\selectfont
\renewcommand{\familydefault}{\sfdefault}
\pagestyle{empty}

\begin{document}
\subsection*{Trigonometría}
\paragraph*{Definición de las funciones trigonométricas} -- \\
\noindent En el siguiente triángulo de lados \scalebox{1.2}{$a$}, \scalebox{1.2}{$b$} y \scalebox{1.2}{$c$} con ángulos \scalebox{1.1}{$A$}, \scalebox{1.1}{$B$} y \scalebox{1.1}{$C$} se definen las siguientes funciones trigonométricas. \\ 

\begin{tikzpicture}
	\coordinate (A) at (0, 0);
	\coordinate (B) at (3.5, 2.1);
	\coordinate (C) at (3.5, 0);
	\draw [thick](B) -- (A) -- (C) -- cycle;

	\draw (-.25, -.15) node{$A$};
	\draw (3.8, 2.25) node{$B$};
	\draw (3.77, -.15) node{$C$};

	\draw (3.75, 1.03) node{$a$};
	\draw (1.8, -.25) node{$b$};
	\draw (1.76, 1.35) node{$c$};

	\draw pic[draw, angle radius=.35cm] {right angle=B--C--A};
\end{tikzpicture}

\noindent En referencia al ángulo $A$:
\begin{itemize}
\item{$\text{seno}\,A = \sen A = \dfrac{a}{c} = \dfrac{\text{Cateto opuesto}}{\text{Hipotenusa}}$}
\item{$\text{coseno}\,A = \cos A = \dfrac{b}{c} = \dfrac{\text{Cateto adyacente}}{\text{Hipotenusa}}$}
\item{$\text{tangente}\,A = \tan A = \dfrac{a}{b} = \dfrac{\text{Cateto opuesto}}{\text{Cateto adyacente}}$}
\item{$\text{cotangente}\,A = \cot A = \dfrac{b}{a} = \dfrac{\text{Cateto adyacente}}{\text{Cateto opuesto}}$}
\item{$\text{secante}\,A = \sec A = \dfrac{c}{b} = \dfrac{\text{Hipotenusa}}{\text{Cateto adyacente}}$}
\item{$\text{cosecante}\,A = \csc A = \dfrac{c}{a} = \dfrac{\text{Hipotenusa}}{\text{Cateto opuesto}}$}
\end{itemize}
\end{document}
